\section{Problem Formulation}
\label{sec:reinforcement_learning_problem_formulation}

This section translates the agent-environment interaction of Figure~\ref{fig:agent_environment_loop} into a mathematical framework.
In this framework, the goal of \gls{rl}, finding a policy that makes sequential decisions so as to maximize the cumulative reward over the long term, can be stated precisely.
The presentation proceeds step by step:
\begin{enumerate}
    \item Modeling the \gls{rl} problem as an \gls{mdp}.
    \item Defining the notion of a policy, which formalizes the agent's behavior.
    \item Describing the interaction between the agent and the environment through trajectories.
    \item Introducing value functions to evaluate policies.
    \item Defining the notion of an optimal policy.
\end{enumerate}

\subsection{Markov Decision Processes}
\label{subsec:markov_decision_processes}

To specify the task and the agent's capacity to perceive and act, the problem is modeled as a \gls{mdp}.
Before providing a more detailed definition, it is useful to clarify the scope of the \gls{mdp} problems considered in this thesis.
We therefore restrict the analysis to the following configuration:

\begin{itemize}
    \item The single-agent learning setting is considered.
    If any other agents are present in the environment, they are treated as part of the environment dynamics.

    \item Finite-horizon \glspl{mdp} are considered.  
    Each interaction sequence between the \gls{mdp} and the policy has a defined beginning and a defined end.  
    Between these two points, the number of interaction steps may vary depending on the situation, but it always remains finite.

    \item The partially observable framework is adopted, where the agent does not have direct access to the true state of the system, but instead relies on limited observations that provide only partial information about the environment.  
\end{itemize}

The resulting framework can be formally described as a single-agent finite-horizon partially observable \gls{mdp}.
For simplicity, we will refer to this setting as an \enquote{\gls{mdp}} throughout this thesis.

A \gls{mdp} is defined by the tuple \((\StateSet, \ObsSet, \ActionSet, \RewardSet, p)\):  
\begin{itemize}
    \item \textbf{\(\StateSet\)}: The set of states, representing all possible configurations of the environment.  
    A state belonging to this set is denoted by \(s \in \StateSet\).  
    Among these, two special subsets are distinguished, the initial states, from which an interaction can begin, and the terminal states, in which the interaction ends.  
    The set of initial states is denoted by \(\StateSetInitial\), with an initial state written as \(\StateInitial \in \StateSetInitial\).
    The set of terminal states is denoted by \(\StateSetTerminal\), with a terminal state written as \(\StateTerminal \in \StateSetTerminal\).
    
    \item \textbf{\(\ObsSet\)}: The set of observations that the agent can perceive from the environment.
    A observation belonging to this set is denoted by \(o \in \ObsSet\). 
    Unlike the state \(s \in \StateSet\), which is hidden from the agent, the observation \(o \in \ObsSet\) provides partial information about the current state.

    \item \textbf{\(\ActionSet\)}: The set of actions that the agent may execute.
    A action belonging to this set is denoted by \(a \in \ActionSet\). 
    Actions influence the dynamics of the \gls{mdp}, thereby shaping its future course.  
    They represent the agent's capacity to act upon its environment.

    \item \textbf{\(\RewardSet\)}: The set of rewards that the agent can receive from the environment.  
    A reward belonging to this set is denoted by \(r \in \RewardSet\). 
    \(\RewardSet\) is a subset of \(\mathbb{R}\).
    Rewards can be either positive or negative, depending on whether the corresponding outcome is considered beneficial or detrimental for the agent.  
    The agent's goal is to maximize the cumulative sum of these rewards over time.

    \item \textbf{\(p\)}: The transition function, which defines the dynamics of the environment.  
    Given a state \(s\) and an action \(a\), this function defines a probability distribution over the joint outcome consisting of the next state \(s'\), the next observation \(o'\), and the reward \(r\).  
    This function is denoted by \(p\), and is defined as:
    \(
    p : \StateSet \times \ActionSet \to \Distributions{\StateSet \times \ObsSet \times \RewardSet}.
    \)
    The notation \((s', o', r) \sim p(\cdot \mid s, a)\) means that the tuple \((s', o', r)\) is sampled from the distribution induced by \(p(\cdot \mid s, a)\).  
    The notation \(p(s', o', r \mid s, a)\) denotes the probability of observing the tuple \((s', o', r)\) under the distribution \(p(\cdot \mid s, a)\).
\end{itemize}

The two quantities, observation \(o\) and state \(s\), are not independent.
They are jointly generated via \((s', o', r) \sim p(\cdot \mid s, a)\), which implicitly defines the relationship between them.
In our \gls{mdp} formalization, this relationship is left implicit because it is never manipulated directly during learning.

\subsection{Policy}
\label{subsec:policy}

The agent's decision-making capability is modeled by a policy function.  
The policy defines the stochastic mapping between observations \(o \in \ObsSet\) and actions \(a \in \ActionSet\).  
A policy function is denoted by \(\pi\), and is defined as:
\(
\pi : \ObsSet \to \Distributions{\ActionSet}.
\)
The notation \(a \sim \pi(\cdot \mid o)\) means that the action \(a\) is sampled from the distribution induced by \(\pi(\cdot \mid o)\).  
The notation \(\pi(a \mid o)\) denotes the probability of selecting action \(a\) under the distribution \(\pi(\cdot \mid o)\).

It is important to note that this policy does not include any explicit mechanism for estimating the underlying state from the observations.
Instead, the policy must learn from experience how to infer the information about the underlying state that is not directly available in the observations.
This reflects the design of many modern \gls{rl} implementations.
Policies are often parameterized by architectures capable of summarizing past observations, such as recurrent networks \parencite{hochreiter1997long, hausknecht2015deep, kapturowski2019recurrent} or attention-based networks \parencite{vaswani2017attention, parisotto2020stabilizing}.
Such a policy maintains an internal state, updated at each time step from the new observation, which acts as a memory of past observations.
For simplicity, this internal state is not made explicit in the notation, and the policy is still written \(\pi(a \mid o)\).
These architectures do not explicitly model the uncertainty over the underlying state.

\subsection{Trajectory}
\label{subsec:trajectory}

Now that both the \gls{mdp} \((\StateSet, \ObsSet, \ActionSet, \RewardSet, p)\) and the policy \(\pi\) have been defined, the interaction between them can be characterized.  
A sequence of consecutive interactions between an environment and a policy is called a trajectory.  
A trajectory is generated through the repeated interaction between the policy and the environment, following the sequence below:
\begin{enumerate}
    \item The environment, being in state \(s\), provides an observation \(o\) to the policy \(\pi\).
    Based on this observation, the policy outputs a probability distribution over actions \(\pi(\cdot \mid o)\), from which an action is sampled:
    \(
    a \sim \pi(\cdot \mid o).
    \)

    \item The sampled action is executed in the environment. The environment then generates a transition:
    \(
    (s', o', r) \sim p(\cdot \mid s, a).
    \)
    It moves to a new state \(s'\), and returns a new observation \(o'\) and a reward \(r\).
\end{enumerate}

Indexing each interaction by a discrete time step \(t\), a trajectory \(\tau\) of length \(T\) can be written as:
\[
\tau = (s_1, o_1, a_1, r_1, s_2, o_2, a_2, r_2, \ldots, s_T, o_T, a_T, r_T),
\]
where \(r_t\) is the reward received after taking action \(a_t\) at time step \(t\).

A trajectory is stochastic, since it is drawn from a distribution determined by the initial state \(s\), the initial observation \(o\), the environment dynamics \(p\), and the policy \(\pi\).
Since the dynamics \(p\) are fixed, they are omitted from the notation.
The probability of observing the trajectory \(\tau\) when following the policy \(\pi\), starting from state \(s\) with observation \(o\), is denoted by \(\Pr_\pi(\tau \mid s, o)\).

An important property of the trajectory distribution is its recursive structure.
Since both the policy \(\pi\) and the transition function \(p\) are known, the distribution over trajectories starting from \(s_1\) with observation \(o_1\) can be expressed in terms of the distribution over trajectories starting from the next state \(s_2\) with observation \(o_2\).
Specifically, the probability of a trajectory \(\tau\) factorizes as:
\[
\Pr_\pi(\tau \mid s_1, o_1) = \pi(a_1 \mid o_1) \; p(s_2, o_2, r_1 \mid s_1, a_1) \; \Pr_\pi(\tau_{2:} \mid s_2, o_2),
\]
where \(\tau_{t:} = (s_t, o_t, a_t, r_t, \ldots, s_T, o_T, a_T, r_T)\) denotes the part of the trajectory starting from time step \(t\).

There exist two special types of trajectories.
First, a trajectory that terminates in a terminal state is called a terminal trajectory, and is denoted by \(\tau^{\mathrm{ter}}\).
Second, a terminal trajectory that starts from an initial state \(\StateInitial\) and an initial observation is called an episode, and is denoted by \(\tau^{\mathrm{ep}}\).

Since trajectories are stochastic, the state, observation, action and reward at each time step are random variables.
Following \textcite{sutton2018reinforcement}, they are denoted by the capital letters \(S_t\), \(O_t\), \(A_t\) and \(R_t\), while the lowercase letters \(s\), \(o\), \(a\) and \(r\) denote particular values that these variables can take.
For instance, \(S_t = s\) expresses the event that the state at time step \(t\) is \(s\).
A trajectory written with lowercase letters, as above, is therefore a realization of this stochastic process.
Expectations over the trajectories generated by the policy \(\pi\) are written \(\mathbb{E}_\pi\).

\subsection{Value Functions}
\label{subsec:value_functions}

To evaluate the ability of a policy \(\pi\) to collect rewards, the notion of return is introduced.

\subsubsection{Return}
\label{subsubsec:return}

The return \(G_t\) is defined as the discounted sum of the rewards collected along a terminal trajectory from time step \(t\) onward:
\[
G_t = R_t + \gamma R_{t+1} + \gamma^2 R_{t+2} + \cdots + \gamma^{T-t} R_T = \sum_{k=t}^{T} \gamma^{k-t} R_k,
\]
where \(\gamma \in [0, 1]\) is the discount factor.
Since the rewards \(R_k\) are random variables, the return \(G_t\) is itself a random variable.

The discount factor \(\gamma\) controls the time horizon of the agent's objective.  
When \(\gamma\) is close to \(0\), the agent focuses almost exclusively on immediate rewards, leading to short-sighted behavior.  
As \(\gamma\) increases toward \(1\), future rewards gain more weight, and the agent adopts a more far-sighted strategy.  
In theory, since the goal is to maximize long-term cumulative reward, the discount factor would ideally be set to \(\gamma = 1\).
In practice, \(\gamma\) is often chosen slightly below \(1\) (e.g., \(0.99\)) to stabilize learning, and to reflect the fact that distant future rewards are inherently more uncertain.

The return also admits a recursive form, which is central to the dynamic programming and \gls{rl} algorithms presented later:
\[
G_t = R_t + \gamma \, G_{t+1}, \qquad G_{T+1} = 0.
\]
For a particular trajectory \(\tau\), the return takes a fixed value, written \(G(\tau) = \sum_{t=1}^{T} \gamma^{t-1} r_t\).
The recursive form then reads \(G(\tau) = r_1 + \gamma \, G(\tau_{2:})\).\footnote{With \(G(\tau_{2:}) = 0\) when \(\tau\) consists of a single interaction step.}

\subsubsection{State-Value Function}
\label{subsubsec:state_value_function}

From the notion of return, functions that estimate the expected return of a policy from a given starting point can be defined.
These functions are called value functions.
The first one is the state-value function, which tells how desirable a state \(s\) with observation \(o\) is under a policy \(\pi\).
Strictly speaking, this function depends on the state \(s\), the observation \(o\), the dynamics \(p\), the discount factor \(\gamma\) and the policy \(\pi\).
Since \(p\) and \(\gamma\) are fixed throughout learning, they are omitted from the notation.
The policy is written as a subscript to emphasize that the value is tied to it, which yields the compact notation \(v_\pi : \StateSet \times \ObsSet \to \mathbb{R}\).

If \(v_\pi(s, o) > v_\pi(s', o')\), then, under policy \(\pi\), being in state \(s\) with observation \(o\) is preferable to being in state \(s'\) with observation \(o'\).
Formally, the state-value function is the expected return when the state is \(s\) and the observation is \(o\), and the policy \(\pi\) is followed from then on:\footnote{By convention, \(v_\pi(s, o) = 0\) for a terminal state \(s \in \StateSetTerminal\).}
\[
v_\pi(s, o) = \mathbb{E}_\pi\big[\, G_t \mid S_t = s, O_t = o \,\big] = \sum_{\tau} \Pr_\pi(\tau \mid s, o) \, G(\tau),
\]
where the sum ranges over the terminal trajectories starting from \((s, o)\).
This value does not depend on the time step \(t\), since neither the dynamics \(p\) nor the policy \(\pi\) change over time.

The state-value function is recursive, a property that follows directly from the recursive form of the return.
Replacing \(G_t\) by \(R_t + \gamma \, G_{t+1}\) in the definition of \(v_\pi\) gives
\[
\begin{aligned}
v_\pi(s, o)
&= \mathbb{E}_\pi\big[\, G_t \mid S_t = s, O_t = o \,\big] \\[4pt]
&= \mathbb{E}_\pi\big[\, R_t + \gamma \, G_{t+1} \mid S_t = s, O_t = o \,\big] \\[4pt]
&= \mathbb{E}_\pi\big[\, R_t + \gamma \, v_\pi(S_{t+1}, O_{t+1}) \mid S_t = s, O_t = o \,\big].
\end{aligned}
\]
The last step holds because, once the next state \(S_{t+1}\) and observation \(O_{t+1}\) are known, the expected value of the remaining return \(G_{t+1}\) is by definition \(v_\pi(S_{t+1}, O_{t+1})\), whatever happened before.

The expectation can then be written explicitly, as in the factorization of \(\Pr_\pi\) given in Subsection~\ref{subsec:trajectory}.
The action \(A_t\) is drawn from \(\pi(\cdot \mid o)\), and the transition \((S_{t+1}, O_{t+1}, R_t)\) that follows is drawn from \(p(\cdot \mid s, A_t)\).
Summing over all their possible values yields the Bellman equation for \(v_\pi\),
\[
\begin{aligned}
v_\pi(s, o)
&= \mathbb{E}_\pi\big[\, R_t + \gamma \, v_\pi(S_{t+1}, O_{t+1}) \mid S_t = s, O_t = o \,\big] \\[4pt]
&= \sum_{a} \pi(a \mid o) \sum_{s', o', r} p(s', o', r \mid s, a) \Big[\, r + \gamma \, v_\pi(s', o') \,\Big].
\end{aligned}
\]
The value of a state is thus expressed through the values of the states that can follow it.

\subsubsection{State-Action Value Function}
\label{subsubsec:state_action_value_function}

The next step is to consider the value of taking a specific action in a given state, in order to compare actions under a given policy \(\pi\).
This function is called the state-action value function \(q_\pi : \StateSet \times \ActionSet \to \mathbb{R}\).
If \(q_\pi(s, a) > q_\pi(s, a')\), then, under policy \(\pi\), choosing action \(a\) in state \(s\) yields a higher expected return than choosing action \(a'\).
It is the expected return when action \(a\) is taken in state \(s\) and the policy \(\pi\) is followed afterwards,
\[
q_\pi(s, a) = \mathbb{E}_\pi\big[\, G_t \mid S_t = s, A_t = a \,\big].
\]
The observation \(O_t\) does not appear in the conditioning, since the observation only influences the choice of the action, which is here fixed to \(a\).

The function \(q_\pi\) can be related to \(v_\pi\) through the same steps as for the Bellman equation.
The only difference is that the first action is no longer drawn from \(\pi\) but fixed to \(a\), which gives
\[
\begin{aligned}
q_\pi(s, a)
&= \mathbb{E}_\pi\big[\, R_t + \gamma \, G_{t+1} \mid S_t = s, A_t = a \,\big] \\[4pt]
&= \mathbb{E}_\pi\big[\, R_t + \gamma \, v_\pi(S_{t+1}, O_{t+1}) \mid S_t = s, A_t = a \,\big] \\[4pt]
&= \sum_{s', o', r} p(s', o', r \mid s, a) \Big[\, r + \gamma \, v_\pi(s', o') \,\Big].
\end{aligned}
\]
The first line uses the recursive form of the return, the second line the definition of \(v_\pi\), and the third line makes the expectation over the transition explicit.

\subsubsection{Relation Between Value Functions}
\label{subsubsec:relation_value_functions}

Each value function can be expressed in terms of the other.
In one direction, the derivation above expresses \(q_\pi\) in terms of \(v_\pi\),
\begin{equation}
q_\pi(s, a) = \sum_{s', o', r} p(s', o', r \mid s, a) \Big[\, r + \gamma \, v_\pi(s', o') \,\Big].
\label{eq:q_from_v}
\end{equation}
In the other direction, the inner sum over \((s', o', r)\) in the Bellman equation for \(v_\pi\) is exactly the right-hand side of Equation~\eqref{eq:q_from_v}.
Replacing it by \(q_\pi(s, a)\) expresses \(v_\pi\) in terms of \(q_\pi\),
\begin{equation}
\begin{aligned}
v_\pi(s, o)
&= \sum_{a} \pi(a \mid o) \underbrace{\sum_{s', o', r} p(s', o', r \mid s, a) \Big[\, r + \gamma \, v_\pi(s', o') \,\Big]}_{q_\pi(s, a)} \\
&= \sum_{a} \pi(a \mid o) \, q_\pi(s, a).
\end{aligned}
\label{eq:v_from_q}
\end{equation}
The value of a state is therefore the average of the values of the actions, weighted by the probability that the policy selects each of them.

\subsection{Optimal Policy}
\label{subsec:optimal_policy}

The goal of \gls{rl} is to find a policy that maximizes the expected return of the episodes it generates.
This policy is called the optimal policy and is denoted by \(\OptimalPolicy\).

Since the policy only receives observations, it must act on an estimate of the underlying state.
The best possible estimate given \(o\) is the distribution \(\Belief{\pi}(s \mid o)\), the probability that the state is \(s\) when \(o\) is observed.
It depends on the policy \(\pi\), since the states in which \(o\) is encountered depend on the actions taken before, and it is learned implicitly by a policy trained on its own experience.
The optimal policy selects the action that maximizes the expected state-action value over this estimate:
\[
\OptimalPolicy(a \mid o) =
\begin{cases}
1 & \text{if } a = \argmax_{a' \in \ActionSet}\;
\sum_{s \in \StateSet} \Belief{\OptimalPolicy}(s \mid o)\, q_{\OptimalPolicy}(s, a'), \\[6pt]
0 & \text{otherwise}.
\end{cases}
\]

As shown, both \(v_\pi\) and \(q_\pi\) can be expressed explicitly in terms of the environment dynamics \(p\) and the policy \(\pi\).
If \(p\) and \(\pi\) are fully known, one can theoretically compute the optimal policy by unfolding the graph of all possible trajectories starting from a given state \(s\) and observation \(o\).
This is possible precisely because \(v_\pi\) and \(q_\pi\) are recursive functions.
Solving an \gls{mdp} in this way is known as dynamic programming.

In practice, most environments have too many states to allow an explicit representation of the dynamics \(p\).  
Storing and computing over the full transition function becomes intractable.
To overcome this limitation, an iterative approach is adopted.
Starting from an initial policy, the policy is progressively improved so that it converges toward \(\OptimalPolicy\).
This iterative improvement through interaction with the environment is precisely what constitutes learning in \gls{rl}, and it is the subject of the next section.